\chapter{Partially coherent images of periodic contact holes}\label{ch:partial}
\paperloc{Illumination, test patterns, and the commercial imaging calculations.}
\reading{Saleh--Teich \S4.4, especially printed pp.~129--134
(PDF pp.~151--156), for coherent transfer functions.
Born--Wolf \S9.5 for coherent/incoherent imaging;
\S10.4 for mutual intensity (printed p.~507 is PDF p.~510);
Hopkins' propagation formula begins at printed p.~512 (PDF p.~505).}

\section{Why one coherent calculation is not the full image}
An extended illumination source contains many emitting regions.
The field arriving from one effective source point is coherent across
the mask for the idealized source-point calculation. Fields from distinct
source points may have rapidly varying relative phases. A resist exposure
averages over times long compared with those phase fluctuations.

Let source point $s$ produce $U_s$ with random phase $\phi_s(t)$.
Then
\[
\avg{\left|\sum_sU_s\ee^{\ii\phi_s(t)}\right|^2}_t
=\sum_s|U_s|^2+
\sum_{s\ne s'}U_sU_{s'}^*
\avg{\ee^{\ii(\phi_s-\phi_{s'})}}_t.
\]
For mutually incoherent source points the last average vanishes.
We must sum \emph{intensities} across source points, even though we
sum \emph{field amplitudes} across diffraction orders from the same point.
Confusing these two sums produces a different physical model.

More generally, the mutual intensity is
\begin{equation}
J(\vct r_1,\vct r_2)
=\avg{U(\vct r_1,t)U^*(\vct r_2,t)}_t,
\qquad
\mu_{12}=\frac{J_{12}}{\sqrt{I_1I_2}}.
\end{equation}
$|\mu_{12}|\leq1$ measures coherence; unity means full coherence at the
two points and zero means no average cross term.
If $U_{\rm out}(x)=\int h(x,r)U_{\rm in}(r)\dd r$, substitution gives
\[
I_{\rm out}(x)=
\iint h(x,r_1)h^*(x,r_2)J_{\rm in}(r_1,r_2)\dd r_1\dd r_2.
\]
The incoherent-source integral below is a convenient special form of this
general propagation rule.

\section{The source-point or Abbe imaging equation}
Let $\vct s$ be an illumination spatial frequency, expressed in
wafer-equivalent coordinates. The incident field on a thin mask is
$\ee^{\ii2\pi\vct s\cdot\vct x}$.
Multiplying the mask function $t$ by this phase shifts its spectrum:
\[
\FT[t(\vct x)\ee^{\ii2\pi\vct s\cdot\vct x}]
=\widetilde t(\vct f-\vct s).
\]
Propagate through pupil $P_z$, change the integration variable to the
mask frequency, and obtain
\begin{equation}
U_{\vct s}(\vct x)=
\ee^{\ii2\pi\vct s\cdot\vct x}
\int\widetilde t(\vct f)P_z(\vct f+\vct s)
\ee^{\ii2\pi\vct f\cdot\vct x}\dd^2f.
\end{equation}
The leading carrier has unit magnitude. For nonnegative source weights
$S$ normalized so that $\int S(\vct s)\dd^2s=1$,
\begin{equation}
I(\vct x;W,z)=
\int S(\vct s)
\left|\int\widetilde t(\vct f)P_z(\vct f+\vct s)
\ee^{\ii2\pi\vct f\cdot\vct x}\dd^2f\right|^2\dd^2s.
\label{eq:abbe}
\end{equation}
The source shifts the spectrum relative to the pupil. It does not
enlarge the physical projection aperture.

\section{The Hopkins transmission cross coefficient}
Expand the squared modulus in Eq.~\eqref{eq:abbe} and interchange
integrals:
\begin{align}
I(\vct x)
={}&\iint\widetilde t(\vct f_1)\widetilde t^*(\vct f_2)
T(\vct f_1,\vct f_2)\\
&\hspace{12mm}\times
\ee^{\ii2\pi(\vct f_1-\vct f_2)\cdot\vct x}
\dd^2f_1\dd^2f_2,\\
T(\vct f_1,\vct f_2)
={}&\int S(\vct s)
P_z(\vct f_1+\vct s)P_z^*(\vct f_2+\vct s)\dd^2s.
\label{eq:tcc}
\end{align}
$T$ is a transmission cross coefficient (TCC). It tells us how strongly
each pair of object frequencies interferes after source averaging.
Its phase contains the pupil-phase difference between the two sampled
locations. The two formulations are equivalent under their shared
thin-mask, scalar, shift-invariant assumptions.

A thick EUV mask can have a source-dependent scattering response.
Then its order amplitudes are $a_{mn}(\vct s,\mathrm{polarization})$,
and one cannot generally factor out a single source-independent
$\widetilde t$ as above. The source-point sum remains useful with the
rigorously computed amplitudes inserted separately.

\section{Diffraction coefficients of a rectangular opening}
Consider an ideal binary opening of width $a$ repeated every $p_x$
and height $b$ repeated every $p_y$. Within a period,
$t(x,y)=1$ inside the centred rectangle and zero elsewhere.
Its Fourier series is
\[
t(x,y)=\sum_{m,n}a_{mn}
\ee^{\ii2\pi(mx/p_x+ny/p_y)}.
\]
Direct integration gives
\begin{align}
a_{mn}
&=\frac{1}{p_xp_y}
\int_{-a/2}^{a/2}\ee^{-\ii2\pi mx/p_x}\dd x
\int_{-b/2}^{b/2}\ee^{-\ii2\pi ny/p_y}\dd y\\
&=\frac{ab}{p_xp_y}
\sinc\!\left(\frac{ma}{p_x}\right)
\sinc\!\left(\frac{nb}{p_y}\right),
\label{eq:rectcoeff}
\end{align}
where $\sinc q=\sin(\pi q)/(\pi q)$ and $\sinc0=1$.
The prefactor is the open-area fraction. A narrower hole spreads its
spectrum more widely, but the pupil still limits the orders that contribute.

For one source point, neglecting the irrelevant carrier,
\begin{equation}
U_{\vct s}(x,y)=\sum_{m,n}a_{mn}
P_z\!\left(\vct s+\left(\frac m{p_x},\frac n{p_y}\right)\right)
\ee^{\ii2\pi(mx/p_x+ny/p_y)}.
\label{eq:orders}
\end{equation}
The integer pair $(m,n)$ names a diffraction order.
Each admitted order is a plane-wave component of the image field.
Their relative phases and amplitudes determine the shape.

\section{Understanding the quasar illumination}
Table 1 specifies inner and outer normalized source radii 0.6 and 0.8,
angle $30^\circ$, and rotation $45^\circ$.
Figure 2 shows four off-axis source regions around diagonal directions.
Their radii are normalized to the relevant projection-pupil scale;
0.8 is not an angle of 0.8 radians.
The exact software interpretation of the angle setting is not documented.
The drawing in this book uses four annular sectors with full angular
width $30^\circ$ as an explicit teaching convention.

At pitch $38\nm$,
\[
\Delta u=\frac{1/p}{f_c}
=\frac{\lambda}{\NA p}
=\frac{13.5}{0.33(38)}\simeq1.07656.
\]
For on-axis illumination the first axial orders lie beyond a unit pupil.
Only the zero order survives in the ideal scalar square lattice,
so it cannot create the desired periodic modulation.
For a source point at radius 0.7 and azimuth $45^\circ$,
$\vct u_s=(0.49497,0.49497)$.
Now the orders $(0,0),(-1,0),(0,-1),(-1,-1)$ all lie inside the pupil.
The oblique source has made several useful orders available simultaneously.

\figteach{quasar_orders}{1.0}{Original pupil diagrams. Left: an idealized
four-sector quasar source. Right: shifted order locations for one diagonal
source point and a 38 nm square pitch. Accepted orders can interfere to
create a two-dimensional contact image. This is a thin-mask explanation,
not a reproduction of the commercial simulator.}

Increasing pitch reduces order spacing in the pupil. Different pitches
therefore sample different phases even when CD stays fixed.
Changing a rectangle's aspect ratio changes the weights $a_{mn}$ in
the two directions. This is why the authors keep many pitches and both
orientations instead of one averaged contact pattern.

\section{A two-order derivation of pattern shift}
Consider
\[
U(x)=a+b\ee^{\ii(qx+\Delta\phi)},\qquad q=2\pi/p,
\]
with real positive $a,b$. Its intensity is
\begin{equation}
I(x)=a^2+b^2+2ab\cos(qx+\Delta\phi).
\label{eq:twoorder}
\end{equation}
The maximum previously at $x=0$ is now at
\begin{equation}
\Delta x=-\frac{\Delta\phi}{q}
=-\frac{p}{2\pi}\Delta\phi.
\label{eq:fringeshift}
\end{equation}
If the differential phase is due to pupil OPD,
$\Delta\phi=2\pi(W_b-W_a)/\lambda$, then
\[
\Delta x=-\frac p\lambda(W_b-W_a).
\]
At $p=38\nm$, a differential OPD of $0.1\nm$ gives a displacement
magnitude of $0.28148\nm$ in this two-order example.
The controlling quantity is a \emph{difference at sampled pupil points},
not the pupil's global RMS.

\section{An even phase perturbation changes shape}
With symmetric side orders,
\[
U(x)=a+b\ee^{\ii\phi}\ee^{\ii qx}
+b\ee^{\ii\phi}\ee^{-\ii qx}
=a+2b\ee^{\ii\phi}\cos qx,
\]
so
\begin{equation}
I(x)=a^2+4b^2\cos^2qx+4ab\cos\phi\cos qx.
\label{eq:threeorder}
\end{equation}
The image remains even in $x$, but its intensity profile changes.
A fixed threshold therefore gives a different printed width.
Because $\cos\phi=1-\phi^2/2+\cdots$, the leading response at this
symmetric operating point is second order.
There need not be a nonzero linear sensitivity to every aberration.

\figteach{phase_to_image}{1.0}{Exact teaching examples. A differential
phase between two orders shifts fringes. A common phase on symmetric
side orders changes an even intensity profile. A real contact combines
many such interference terms and an incoherent source average.}

\section{Defocus, symmetry, and source design}
An order at normalized radius $\rho$ acquires the defocus phase
$-\pi z\NA^2\rho^2/\lambda$ in the paraxial model.
Two orders at equal radius receive the same defocus phase and preserve
their relative phase. This is one reason suitable off-axis illumination
can improve focus robustness for selected patterns.
For unequal radii, the difference is
\[
\Delta\phi_z=-\frac{\pi z\NA^2}{\lambda}
(\rho_2^2-\rho_1^2).
\]
One source design cannot generally equalize every important order pair
for all pitches and both rectangular dimensions.
Source optimization is a pattern-dependent interference problem.

\chapter{From light intensity to CD, PS, PVB, and NILS}\label{ch:metrics}
\paperloc{The six imaging indicators used throughout the paper.}

\section{An aerial image is not yet a resist contour}
The aerial image is the optical intensity in the image plane.
Exposure dose is energy per area, obtained by integrating intensity
over time. A full resist model may include photon absorption,
secondary electrons, chemical reactions, diffusion, and development.
The paper does not provide enough resist-model detail to reconstruct
that stage exactly.

To derive the basic geometry, use a declared constant-threshold model:
\begin{equation}
D\,I(x,y;z,W)=T.
\label{eq:threshold}
\end{equation}
$D$ is a relative dose multiplier and $T$ is a fixed exposure threshold
in consistent units. This is a teaching approximation. It captures how
intensity perturbations move edges without claiming a full EUV resist model.
For a bright contact opening, the printed interior is taken to be the
region above threshold; resist tone can reverse that interpretation.

\section{Critical dimension and centre position}
On a specified horizontal cut through a feature, let threshold crossings
be $x_L$ and $x_R$. Define
\begin{equation}
\mathrm{CD}_x=x_R-x_L,\qquad
x_c=\frac{x_R+x_L}{2},\qquad
\mathrm{PS}_x=x_c-x_{\rm target}.
\label{eq:CDPS}
\end{equation}
Use analogous definitions in $y$. For a nonrectangular contour, the
measurement cut, bounding box, fitted shape, or centroid definition must
be specified. A centroid computed from the whole intensity distribution
need not equal the midpoint between two threshold edges.

Small edge changes satisfy
\begin{equation}
\delta\mathrm{CD}_x=\delta x_R-\delta x_L,\qquad
\delta\mathrm{PS}_x=\frac{\delta x_R+\delta x_L}{2}.
\label{eq:edgecombine}
\end{equation}
Equal edge motion shifts the pattern. Opposite edge motion changes its
width. These two combinations explain why CD and PS provide different
information about a wavefront.

\section{Deriving the edge-displacement formula}
Let an original edge satisfy $I_0(x_e)=T$ at nominal dose one.
Perturb the image by $\delta I$ and the edge by $\delta x$:
\[
I_0(x_e+\delta x)+\delta I(x_e+\delta x)=T.
\]
Keeping first-order terms,
$I_0(x_e)+I_0'(x_e)\delta x+\delta I(x_e)=T$.
Cancel the original threshold equation:
\begin{equation}
\delta x=-\frac{\delta I(x_e)}{I_0'(x_e)}.
\label{eq:edge}
\end{equation}
The same intensity error causes a larger edge displacement when the
edge slope is shallow.
For a two-dimensional contour, normal displacement is
\[
\delta n=-\frac{\delta I}{|\nabla I_0|},
\]
with the signed normal chosen along increasing intensity.
At a vanishing slope, the linear formula fails: an opening may appear,
disappear, split, or merge. These topology changes also make learned
metric maps difficult to approximate smoothly.

\section{NILS makes the slope dimensionless}
For characteristic width $w$, define the magnitude of normalized image
log slope at an edge:
\begin{equation}
\mathrm{NILS}=w\left|\frac{\partial\ln I}{\partial x}\right|_{x_e}
=w\frac{|I'(x_e)|}{I(x_e)}.
\label{eq:nils}
\end{equation}
Different reporting conventions retain or remove the slope sign.
The dimensionless magnitude is what matters for this sensitivity argument.

A relative dose change $\epsilon=\delta D/D$ multiplies the image
by approximately $1+\epsilon$, so $\delta I=\epsilon I$.
Equation \eqref{eq:edge} gives
\begin{equation}
|\delta x|\simeq\frac{w}{\mathrm{NILS}}|\epsilon|.
\label{eq:nilsdose}
\end{equation}
For a symmetric bright feature with equal slope magnitudes at the two
edges, both edges move outward when dose increases:
\[
|\delta\mathrm{CD}|\simeq
\frac{2w}{\mathrm{NILS}}|\epsilon|.
\]
The factor of two comes from two moving edges. It must not be silently
carried into a one-edge placement formula.

\worked{A 5 percent dose perturbation}
If $w=14\nm$ and NILS is 2, one edge moves approximately
$14(0.05)/2=0.35\nm$ under a 5 percent dose change.
The full CD changes by approximately $0.70\nm$ in the symmetric example.
These assumed slopes are not measured values from the paper.
The example explains why contrast loss can enlarge process variation
even if nominal CD is adjusted back to target by changing dose.

\section{Focus-exposure matrices and process windows}
For each focus $z$ and relative dose $D$, simulate a contour and compute
$\mathrm{CD}(z,D)$, $\mathrm{PS}(z,D)$, and other required measures.
Plots of CD against focus at several doses are often called Bossung curves.
A process window is the set of focus-dose pairs satisfying \emph{all}
chosen acceptance conditions:
\[
\mathcal P=\{(z,D):\mathrm{CD}_{\min}\leq\mathrm{CD}(z,D)
\leq\mathrm{CD}_{\max},\ |\mathrm{PS}(z,D)|\leq b,\ldots\}.
\]
An attractive nominal image at $z=0,D=1$ can coexist with a very small
process window. Aberration-induced sensitivity to focus and dose is
therefore physically significant.

The paper gives $z=-15,0,+15\nm$ and dose variation $\pm5\%$.
If one uses doses $0.95,1.00,1.05$, a teaching focus-dose matrix has
nine conditions. This is a declared sampling choice here, not proof that
the authors used exactly that grid or only those dose samples.

\section{What a process variation band can mean}
For contours $C(z,D)$, their envelope is the region swept out by edge
positions over the chosen process conditions.
On a horizontal cut, define edge-band widths
\begin{equation}
B_L=\max x_L-\min x_L,\qquad
B_R=\max x_R-\min x_R.
\label{eq:pvbedges}
\end{equation}
A scalar horizontal PVB might be the larger edge-band width, a
specified average, or another software-defined contour measure.
By contrast, CD range is
\begin{equation}
B_{\rm CD}=\max(x_R-x_L)-\min(x_R-x_L).
\label{eq:cdspan}
\end{equation}
These quantities are not generally equal. For pure translation,
$B_{\rm CD}=0$ while both edge bands can be large.
For symmetric expansion/contraction, each edge band can be half
the CD range.

\careful{The paper does not give a mathematical extraction rule for
$\mathrm{PVB}_x$ and $\mathrm{PVB}_y$. Use its term as a reported
process-contour variation indicator. Do not assume without verification
that it is simply maximum CD minus minimum CD. The laboratory in this
book reports the specific edge-band widths of Eq.~\eqref{eq:pvbedges}.}

\section{Placement error, overlay, and edge placement error}
Pattern shift concerns one feature relative to its intended location.
Overlay concerns alignment between layers or exposures.
An aberration-induced shift can consume part of an overlay allowance,
but stage alignment, mask registration, wafer distortion, and other
contributors also matter.

For a one-dimensional target with left and right edges, small
edge placement errors combine CD and PS:
\begin{equation}
\mathrm{EPE}_R=\mathrm{PS}+\frac{\delta\mathrm{CD}}2,\qquad
\mathrm{EPE}_L=\mathrm{PS}-\frac{\delta\mathrm{CD}}2.
\label{eq:epe}
\end{equation}
This algebra shows why controlling only the feature centre is insufficient
to control every edge. A contact can be correctly centred but too large.
Conversely, its width can be perfect while its entire position is wrong.

\section{Why horizontal and vertical errors differ}
The mask incidence plane establishes a preferred direction. A rectangular
hole has different spectra in $x$ and $y$. A quasar source has discrete
angular structure. A wavefront mode has its own orientation.
The interplay can make one direction more sensitive than the other.

The paper's relatively large $\mathrm{PS}_x$ and small
$\mathrm{PS}_y$ should not be interpreted as a universal property of all
EUV contacts. The incidence azimuth, coordinate conventions, baseline
placement, and definition of the reported shifts would need to be known.
A small denominator in a percentage error is a separate statistical issue,
addressed in Chapter \ref{ch:errors}.

\chapter{How an aberration coefficient changes a metric}\label{ch:sensitivity}
\section{Differentiating the pupil and the field}
Let $W=\sum_jc_jZ_j$ with coefficients in length units.
Then
\begin{equation}
\frac{\partial P}{\partial c_j}
=\frac{\ii2\pi}{\lambda}Z_jP.
\label{eq:pupilDerivative}
\end{equation}
For a fixed thin-mask spectrum and one source point,
\[
\frac{\partial U_{\vct s}}{\partial c_j}
=\frac{\ii2\pi}{\lambda}
\int \widetilde t(\vct f)\,
Z_j\!\left(\frac{\vct f+\vct s}{f_c}\right)
P(\vct f+\vct s)\ee^{\ii2\pi\vct f\cdot\vct x}\dd^2f.
\]
Differentiate $I=UU^*$:
\begin{equation}
\frac{\partial I}{\partial c_j}
=2\Real\left\{U^*\frac{\partial U}{\partial c_j}\right\}.
\label{eq:intDerivative}
\end{equation}
For a partially coherent source, integrate this expression over source
weights. In a vector calculation, replace the scalar product with the
appropriate polarization sum.

The derivative is physically a measure of how a phase change in one
mode perturbs the constructive and destructive interference in the image.
It is not determined merely by how large the mode looks in a pupil plot.

\section{Differentiating a printed edge}
Apply implicit differentiation to
$I(x_e(\vct c),\vct c)=T$:
\[
\frac{\partial I}{\partial x}\frac{\partial x_e}{\partial c_j}
+\frac{\partial I}{\partial c_j}=0.
\]
Therefore
\begin{equation}
\frac{\partial x_e}{\partial c_j}
=-\frac{\partial I/\partial c_j}{\partial I/\partial x}.
\label{eq:edgeSensitivity}
\end{equation}
Combine left and right derivatives using Eq.~\eqref{eq:edgecombine}
to obtain the sensitivity of CD and PS. This supplies an explicit
physics-based path from coefficient to metric.

For $M$ measured indicators, define
\begin{equation}
J_{ij}=\left.\frac{\partial F_i}{\partial c_j}\right|_{\vct c_0},
\qquad
\delta\vct y\simeq J\,\delta\vct c.
\label{eq:sensitivity}
\end{equation}
$J$ is the imaging Jacobian. Its rows correspond to particular patterns,
orientations, and metrics; its columns correspond to modes.
If coefficients and metrics are both in nm, its entries are dimensionless.
If coefficients are in milliwaves, the entries have units
nm per milliwave.

\section{A numerical derivative when an analytic derivative is unavailable}
With a black-box simulator, a centred finite difference gives
\begin{equation}
J_{ij}\simeq
\frac{F_i(\vct c_0+h\vct e_j)-F_i(\vct c_0-h\vct e_j)}{2h}.
\label{eq:finitedifference}
\end{equation}
Here $\vct e_j$ changes only coefficient $j$.
Taylor expansion shows the leading truncation error is $O(h^2)$ for
a sufficiently smooth map.
Making $h$ extremely small is not always beneficial: numerical contour
interpolation, solver convergence error, and cancellation can dominate.
Check derivative stability over a sensible range of $h$ and ensure
the same physical feature is being measured.

\section{Interactions and nonlinear response}
The second-order expansion of indicator $i$ is
\begin{equation}
\delta y_i\simeq\sum_jJ_{ij}\delta c_j+
\frac12\sum_{j,k}H_{ijk}\delta c_j\delta c_k,
\end{equation}
where $H_{ijk}$ is a second derivative. Terms with $j\ne k$ describe
interactions between modes. Single-mode sweeps do not by themselves
measure those cross terms.

For example, if $y=c_1c_2$, varying either coefficient while keeping the
other zero produces no response. Yet both nonzero together produce a
finite response. This elementary example explains why the paper
supplements individual-mode sweeps with combined aberrations.

\section{Symmetry predicts some zero sensitivities}
If a complete imaging configuration is invariant under reflection
$x\mapsto-x$, an even phase perturbation cannot by itself generate an
odd placement change under conditions preserving that symmetry.
Likewise, a quantity may depend on a coefficient through $c^2$ near zero,
giving zero first derivative but nonzero second derivative.

These conclusions require symmetry of the whole experiment: source,
mask, pupil amplitude, aberration, focus, and metric definition.
Oblique EUV mask scattering can destroy a symmetry present in a
simplified scalar calculation. Mode names alone cannot establish
which physical direction shifts.

\section{Why a physical inverse is difficult but not impossible}
The paper says required aberrations cannot be deduced using physical
laws or equations because of complexity. The careful interpretation is
that a convenient closed-form inverse is generally unavailable.
The forward map is still governed by wave physics; its derivatives can
be computed, and iterative constrained optimization can use them.
The neural network offers an approximate, data-driven route to candidate
solutions. It does not replace Maxwell's equations or prove that
physics-based inverse design is impossible.

\checkpoint{Can a coefficient with a small derivative at zero be given
an unlimited budget? No. Higher-order effects, interactions, changes in
the active contour, and other patterns can become important away from zero.}
